\documentclass[10pt]{article}

% Specify the margins and text size.
\setlength{\textwidth}{6.4in}
\setlength{\textheight}{9.5in}
\setlength{\oddsidemargin}{0pt}
\setlength{\evensidemargin}{0pt}
\setlength{\topmargin}{0pt}
\setlength{\hoffset}{.05in}
\setlength{\voffset}{-1in}

\setlength{\parskip}{5pt}
\setlength{\parindent}{0pt}

% Load some fonts and symbol packages
\usepackage{latexsym}
\usepackage{pifont}       % contains 'star' symbol for counterinsurgency handbook title
\usepackage{yfonts} 
\usepackage{amsmath}
\usepackage{amsfonts}

\usepackage{graphicx}     % actually, this is loaded by pstricks
\usepackage[T1]{fontenc}
\usepackage{ifthen}
\usepackage{pstricks,pst-grad,pst-text,pst-node,multido,pst-plot,pst-func,calc,pst-3dplot}
\usepackage[all]{xy}
%\usepackage{animate}

% The hyperref package inserts links into the pdf.
\definecolor{MyLinkColor}{rgb}{.1,.2,1}
\definecolor{MyCiteColor}{rgb}{.1,1,.2}
\definecolor{MyURLColor}{rgb}{.4,.4,.4}
\usepackage[backref=true,pagebackref=false,hyperindex,colorlinks=true,
  linkcolor=MyLinkColor,urlcolor=MyURLColor]{hyperref}


% The tweaklist package is something I found on the web.  It provides a simple interface
% for making changes to spacing used in the itemize and enumerate environments.  Comment
% this out if you don't care to use tweaklist.
\usepackage{tweaklist}
\renewcommand{\itemhook}{\setlength{\parskip}{2pt}\setlength{\parsep}%
{1pt}\setlength{\topsep}{0pt}\setlength{\itemsep}{0pt}}

\newcommand{\U}{\underline{\hspace{5pt}}}

\usepackage{listings}
\newcommand{\Z}{\hphantom{0}}

\begin{document}
\pagestyle{empty}
\lstset{language=R, showspaces=false, showstringspaces=false}

\href{http://www.su.edu}{\includegraphics[height=1.75cm]{sulogo.eps}}
\vspace{-1.69cm}

{{\ }\hfill\small
\begin{tabular}{cl}
& Math 207\\
& Introduction to Statistics\\
\end{tabular}
}
\setlength{\baselineskip}{1.05\baselineskip}

\begin{center}
\textbf{\large  Exam IV}
\end{center}
%No books, notes or electronic devices may be used on this exam.
\medskip

1. (25 points) A survey is conducted of 100 freshmen at public universities.
It is found that these students work 11.5 hours a week for pay, 
on average.  The SD of the data is 5.  

\hspace{20pt} a) Find a 95\% confidence interval
on the  hours  a week worked by freshmen at public universities.
\vspace{2.5in}

\hspace{20pt} b) True or false and explain:  There is a 95\% chance that
if a second sample of 100 freshmen were taken
the average of the sample would be in the range you found in a).
\vspace{1.5in}

\hspace{20pt} c) True or false and explain:  There is a 95\% chance that the
average hours a week worked for pay among \textit{all} freshmen at public universities
is in the range you found in a).
\vspace{1.5in}

\hspace{20pt} d) True or false and explain:  There is a 95\% chance that the
average hours a week worked for pay among all 18--20 year olds
is in the range you found in a).
\vspace{1in}
\vfill
\eject
{\ }



2. (15 points)
Would taking the average of 25 measurements divide the likely size of the chance error by a factor of 5, 10 or 25? Justify your answer.
\vspace{2.5in}

3. (15 points) A surveyor is measuring the distances between five points $A$, $B$, $C$,
$D$, and $E$ along a straight line.  She finds that each of the four distances measures
one mile, give or take an inch or so.  The distance from $A$ to $E$ is about four miles; 
but the estimate is likely to be off by around:  
4~inches, 2~inches, 1~inch, 1/2~inch or 1/4~inch. Justify your answer.
\vfill

4. (20 points) Five hundred draws are made at random from the box
\begin{center}
\begin{pspicture}(0,0.1)(4,0.3)
\psline(0,0.5)(0,0)(4,0)(4,0.5)
\rput[l](0.15,0.3){60,000\ \framebox{0}\,s}
\rput[l](2.15,0.3){20,000\ \framebox{1}\,s}
\end{pspicture}
\end{center}
Note that $\sqrt{0.25\times 0.75}\,/\sqrt{500} \approx 0.02$. True or false? 
\smallskip

\begin{itemize}
\item The expected value for the percentage of 1s among the draws is exactly 25\%.\\
\item The expected value for the percentage of 1s among the draws is around 25\%, give
or take 2\% or so.\\
\item The percentage of 1s among the draws will be around 25\%, give or take 2\% or so.\\
\item The percentage of 1s among the draws will be exactly 25\%.\\
\item The percentage of 1s in the box is exactly 25\%.\\
\item The percentage of 1s in the box is around 25\%, give or take 2\% or so.
\end{itemize}

\eject
{\ }

5.  (20 points) Based on the survey data in \#1, a news organization reports that the
freshmen at public universities work~12 hours a week for pay, on average.
Use a $z$-test to comment on this assertion.
State the null and alternative hypotheses. Calculate the test statistic and 
estimate the $P$ value.
%z-test
\vspace{6in}

6. (10 points)  
Fill in the blanks: 
\bigskip

The \underline{\hspace{60pt}} hypothesis says that the difference is due to chance
but the \underline{\hspace{60pt}} says that the difference is real.  
\vfill
\eject
%%%%%%%%%%%%%%%%%%%%%%%%%%%%%%%%%%%%%%%%%%%%%%%%%%%%%%%%%%%%%%%%%%%%%%%%%%%

\href{http://www.su.edu}{\includegraphics[height=1.75cm]{sulogo.eps}}
\vspace{-1.69cm}

{{\ }\hfill\small
\begin{tabular}{cl}
& Math 207\\
& Introduction to Statistics\\
\end{tabular}
}
\begin{center}
\textbf{\large  Exam IV}
\end{center}
%No books, notes or electronic devices may be used on this exam.
\medskip

1. (25 points) A survey is conducted of 100 freshmen at public universities.
Among the surveyed  students, 20\% (or 1/5) work at least 15 hours per week for pay.

\hspace{20pt} a) Find a 95\% confidence interval
on the  percent of freshmen at public universities who work at least 15 hours per week for pay.
\vspace{2.3in}

\hspace{20pt} b) True or false and explain:  There is a 95\% chance that the
percent among \textit{all} freshmen at public universities is in the range you found in a).
\vspace{1.5in}

\hspace{20pt} c) True or false and explain:  There is a 95\% chance that
if a second sample of 100 freshmen were taken,
the sample percent would be in the range you found in a).
\vspace{1.5in}

\hspace{20pt} d) True or false and explain:  There is a 95\% chance that the
average hours a week worked for pay among all 18--20 year olds
is in the range you found in a).
\vspace{1in}
\vfill
\eject
{\ }



2. (10 points)
Would taking the sum of 25 measurements multiply
the likely size of the chance error by a factor of 5, 10 or 25? Justify your answer.
\vspace{1.8in}

3. (10 points) A surveyor is measuring the distances between five points $A$, $B$, $C$,
$D$, and $E$ along a straight line.  She finds that each of the four distances measures one 
mile, give or take an inch or so.  The distance from $A$ to $E$ is about four miles; 
but the estimate is likely to be off by around:  
4~inches, 2~inches, 1~inch, 1/2~inch or 1/4~inch. Justify your answer.
\[\xymatrix{
A\ar@{-}[r] & B\ar@{-}[r] & C\ar@{-}[r] & D\ar@{-}[r] & E}\]
\vspace{1.8in}

4.  (15 points) The Zorro News Organization reports that  50\%
of freshmen at public universities  work 15 hours a week for pay.
If that is report is accurate, approximately what is the chance of getting a sample
percent as low as 20\% (as we did in \#1) or lower?  Hint:  Use a $z$-test.
\vfill
\eject


5. (20 points) Five hundred draws are made at random from the box
\begin{center}
\begin{pspicture}(0,0.1)(4,0.3)
\psline(0,0.5)(0,0)(4,0)(4,0.5)
\rput[l](0.15,0.3){60,000\ \framebox{0}\,s}
\rput[l](2.15,0.3){20,000\ \framebox{1}\,s}
\end{pspicture}
\end{center}
Note that $\sqrt{0.25\times 0.75}\,/\sqrt{500} \approx 0.02$. True or false?  Explain your conclusions.
\smallskip

\begin{itemize}
\item[i.] The expected value for the percentage of 1s among the draws is exactly 25\%.\\[1in]
\item[ii.] The expected value for the percentage of 1s among the draws is around 25\%, give
or take 2\% or so.\\[1in]
\item[iii.] The percentage of 1s among the draws will be around 25\%, give or take 2\% or so.\\[1in]
\item[iv.] The percentage of 1s among the draws will be exactly 25\%.\\[1in]
\item[v.] The percentage of 1s in the box is exactly 25\%.\\[1in]
\item[vi.] The percentage of 1s in the box is around 25\%, give or take 2\% or so.
\end{itemize}

\vfill
\eject

6. (20 points) If a  particular scientific instrument is working correctly, it's averge measurement
should be~9.  It is tested twice and the results are 9 and 11.  Run a $t$ test to see if the
difference between the observed and expected averages is significant.  State the null and alternative hypotheses.


\vfill
\begin{tabular}{ccc}
\begin{tabular}{|c|cccc|}\hline
Degrees\vphantom{\LARGE Y} & &&&\\
of      & &\rput(20pt, 10pt){$t$}&&\\
Freedom & $-2.0$ & $-1.5$ & $-1.0$ & $-0.5$\\[3pt]\hline
1       & $0.148$ & $0.187$ & $0.250$ & $0.352$\vphantom{\LARGE Y}\\
2       & $0.092$ & $0.136$ & $0.211$ & $0.333$\\
3       & $0.070$ & $0.115$ & $0.196$ & $0.326$\\\hline
\end{tabular}
&\hspace{1in} &
\definecolor{lg}{rgb}{.6,.6,.6}
\begin{pspicture}(-2.5,0.5)(2.5,1.5)
\psframe[linewidth=0.02](-3.2,-0.88)(3.2,2.05)
\rput(0,1.7){$t$ distribution with $2$ degrees of freedom}
\psset{yunit=4}
\pscustom[linestyle=none,fillstyle=solid,fillcolor=lg]{%
\psTDist*[linewidth=0,linecolor=blue,nue=2]{-2.25}{-0.5}
\psline(-0.5,0.2963)(-0.5,0)(-2.25,0)(-2.25,0.0533)
}
\psTDist[linecolor=blue,nue=2]{-2.25}{2.25}
\psline{<->}(-2.5,0)(2.5,0)
\psline(-0.5,0)(-0.5,-0.05)\rput[t](-0.5,-0.08){\footnotesize$-0.5\hphantom{-}$}
\psline(-2,0)(-2,-0.05)    \rput[t](-2,-0.08){\footnotesize$-2\hphantom{-}$}
\psline(-1,0)(-1,-0.05)   % \rput[tr](-1,-0.08){\footnotesize$-1$}
\psline(0,0)(0,-0.05)      \rput[t]( 0,-0.08){\footnotesize$0$}
\psline(1,0)(1,-0.05)      \rput[t]( 1,-0.08){\footnotesize$1$}
\psline(2,0)(2,-0.05)      \rput[t]( 2,-0.08){\footnotesize$2$}
\psbezier{->}(-6.13,-0.05)(-4,0.3)(-2,0.5)(-1,0.1)
\end{pspicture}
\end{tabular}

\eject

\begin{center}
\textbf{\large Formula Sheet}
\end{center}\vspace{-15pt}
\begin{align*}
\mbox{mean} &=\frac{1}{n}\left(x_1 + \cdots + x_n\right)\\[5pt]
\mbox{SD}   &=\sqrt{\frac{1}{n}\left( (x_1 - \mbox{mean})^2 + \cdots + (x_n - \mbox{mean})^2\right)}\\[10pt]
\mbox{sd}   &=\sqrt{\frac{1}{n-1}\left( (x_1 - \mbox{mean})^2 + \cdots + (x_n - \mbox{mean})^2\right)}\\[10pt]
%
\mbox{EV}_{\mbox{\footnotesize sum}} & = n\,\mbox{AV}_{\mbox{\footnotesize box}}\\
\mbox{SE}_{\mbox{\footnotesize sum}} & = \sqrt{n}\;\;\mbox{SD}_{\mbox{\footnotesize box}}\\[10pt]
%
\mbox{EV}_{\mbox{\footnotesize av}} & = \mbox{AV}_{\mbox{\footnotesize box}}\\
\mbox{SE}_{\mbox{\footnotesize av}} & = \mbox{SD}_{\mbox{\footnotesize box}}\,\big{/}\sqrt{n}\\[10pt]
%
%\mbox{EV}_{\mbox{\footnotesize \%}} & = \mbox{AV}_{\mbox{\footnotesize box}}\\
%\mbox{SE}_{\mbox{\footnotesize \%}} & = \mbox{SD}_{\mbox{\footnotesize box}}\,\big{/}\sqrt{n}\\[5pt]
%
\noalign{Shortcut formula (if there are only two different kinds of numbers in the box):}
\mbox{SD}_{\mbox{\footnotesize box}} &=
  \left(\mbox{big \#} - \mbox{small \#}\right)\,\sqrt{
  \left(\begin{array}{c}\mbox{fraction of}\\ \mbox{tickets with the}\\ \mbox{big number}\end{array}\right)
  \left(\begin{array}{c}\mbox{fraction of}\\\mbox{tickets with the}\\\mbox{small number}\end{array}\right)}\\
%
\mbox{95\% confidence interval} &= \mbox{observed} \;\pm\; 2\,\mbox{SE}\\[5pt]
z &= \frac{\mbox{observed} - \mbox{expected}}{\mbox{standard error}}
\end{align*}
\vfill

\begin{center}
\textbf{\large Formula Sheet}
\end{center}\vspace{-15pt}
\begin{align*}
\mbox{mean} &=\frac{1}{n}\left(x_1 + \cdots + x_n\right)\\[5pt]
\mbox{SD}   &=\sqrt{\frac{1}{n}\left( (x_1 - \mbox{mean})^2 + \cdots + (x_n - \mbox{mean})^2\right)}\\[10pt]
\mbox{sd}   &=\sqrt{\frac{1}{n-1}\left( (x_1 - \mbox{mean})^2 + \cdots + (x_n - \mbox{mean})^2\right)}\\[5pt]
%
\mbox{EV}_{\mbox{\footnotesize sum}} & = n\,\mbox{AV}_{\mbox{\footnotesize box}}\\
\mbox{SE}_{\mbox{\footnotesize sum}} & = \sqrt{n}\;\;\mbox{SD}_{\mbox{\footnotesize box}}\\[10pt]
%
\mbox{EV}_{\mbox{\footnotesize av}} & = \mbox{AV}_{\mbox{\footnotesize box}}\\
\mbox{SE}_{\mbox{\footnotesize av}} & = \mbox{SD}_{\mbox{\footnotesize box}}\,\big{/}\sqrt{n}\\[10pt]
%
%\mbox{EV}_{\mbox{\footnotesize \%}} & = \mbox{AV}_{\mbox{\footnotesize box}}\\
%\mbox{SE}_{\mbox{\footnotesize \%}} & = \mbox{SD}_{\mbox{\footnotesize box}}\,\big{/}\sqrt{n}\\[5pt]
%
\noalign{Shortcut formula (if there are only two different kinds of numbers in the box):}
\mbox{SD}_{\mbox{\footnotesize box}} &=
  \left(\mbox{big \#} - \mbox{small \#}\right)\,\sqrt{
  \left(\begin{array}{c}\mbox{fraction of}\\ \mbox{tickets with the}\\ \mbox{big number}\end{array}\right)
  \left(\begin{array}{c}\mbox{fraction of}\\\mbox{tickets with the}\\\mbox{small number}\end{array}\right)}\\[10pt]
%
\mbox{95\% confidence interval} &= \mbox{observed} \;\pm\; 2\,\mbox{SE}\\[5pt]
z &= \frac{\mbox{observed} - \mbox{expected}}{\mbox{standard error}}
\end{align*}
\eject

%%%%%%%%%%%%%%%%%%%%%%%%%%%%%%%%%%%%%%%%%%%%%%%%%%%%%%%%%%%%%%%%%%%%%%%%%%%

\href{http://www.su.edu}{\includegraphics[height=1.75cm]{sulogo.eps}}
\vspace{-1.69cm}

{{\ }\hfill\small
\begin{tabular}{cl}
& Math 207\\
& Introduction to Statistics\\
\end{tabular}
}
\begin{center}
\textbf{\large  Exam IV}
\end{center}
%No books, notes or electronic devices may be used on this exam.
\medskip

1. (25 points) A survey is conducted of 100 freshmen at public universities.
Among the surveyed  students, 20\% (or 1/5) work at least 15 hours per week for pay.

\hspace{20pt} a) Find a 95\% confidence interval
on the  percent of freshmen at public universities who work at least 15 hours per week for pay.
\vspace{2.3in}

\hspace{20pt} b) True or false and explain:  There is a 95\% chance that the
percent among \textit{all} freshmen at public universities is in the range you found in a).
\vspace{1.5in}

\hspace{20pt} c) True or false and explain:  There is a 95\% chance that
if a second sample of 100 freshmen were taken,
the sample percent would be in the range you found in a).
\vspace{1.5in}


\hspace{20pt} d) True or false and explain:  There is a 95\% chance that 
 among all 18--20 year olds, the percent who work at least 15 hours per week for pay 
is in the range you found in a).
\vspace{1in}
\vfill
\eject
{\ }



2. (15 points)
Would taking the average of 25 measurements decrease
the likely size of the chance error by a factor of 5, 10 or 25? Justify your answer.
\vspace{1.5in}

3. (15 points) A surveyor is measuring the distances between five points $A$, $B$, $C$,
$D$, and $E$ along a straight line.  She finds that each of the four distances measures one 
mile, give or take an inch or so.  The distance from $A$ to $E$ is about four miles; 
but the estimate is likely to be off by around:  
4~inches, 2~inches, 1~inch, 1/2~inch or 1/4~inch. Justify your answer.\vspace{-5pt}

\[\xymatrix{
A\ar@{-}[r] & B\ar@{-}[r] & C\ar@{-}[r] & D\ar@{-}[r] & E}\]
\vspace{2in}


4.  (20 points) The Zorro News Organization reports that  50\%
of freshmen at public universities  work 15 hours a week for pay.
If that is report is accurate, approximately what is the chance of getting a sample
percent as low as or lower than 20\% (as we did in \#1)?  Hint:  calculate $z$ then a $p$-value (area under the normal curve).
\vfill
\eject


5. (25 points) Five hundred draws are made at random from the box
\begin{center}
\begin{pspicture}(0,0.1)(4,0.3)
\psline(0,0.5)(0,0)(4,0)(4,0.5)
\rput[l](0.15,0.3){60,000\ \framebox{0}\,s}
\rput[l](2.15,0.3){20,000\ \framebox{1}\,s}
\end{pspicture}
\end{center}
Note that $\sqrt{0.25\times 0.75}\,/\sqrt{500} \approx 0.02$. True or false?  Explain your conclusions.
\smallskip

\begin{itemize}
\item[i.] The expected value for the percentage of 1s among the draws is exactly 25\%.\\[1in]
\item[ii.] The expected value for the percentage of 1s among the draws is around 25\%, give
or take 2\% or so.\\[1in]
\item[iii.] The percentage of 1s among the draws will be around 25\%, give or take 2\% or so.\\[1in]
\item[iv.] The percentage of 1s among the draws will be exactly 25\%.\\[1in]
\item[v.] The percentage of 1s in the box is exactly 25\%.\\[1in]
\item[vi.] The percentage of 1s in the box is around 25\%, give or take 2\% or so.
\end{itemize}

\vfill
\eject

\begin{center}
\textbf{\large Formula Sheet}
\end{center}


\begin{align*}
\mbox{mean} &=\frac{1}{n}\left(x_1 + \cdots + x_n\right)\\[10pt]
\mbox{SD}   &=\sqrt{\frac{1}{n}\left( (x_1 - \mbox{mean})^2 + \cdots + (x_n - \mbox{mean})^2\right)}\\[10pt]
\mbox{sd}   &=\sqrt{\frac{1}{n-1}\left( (x_1 - \mbox{mean})^2 + \cdots + (x_n - \mbox{mean})^2\right)}\\[10pt]
%
\mbox{EV}_{\mbox{\footnotesize sum}} & = n\,\mbox{AV}_{\mbox{\footnotesize box}}\\
\mbox{SE}_{\mbox{\footnotesize sum}} & = \sqrt{n}\;\;\mbox{SD}_{\mbox{\footnotesize box}}\\[10pt]
%
\mbox{EV}_{\mbox{\footnotesize av}} & = \mbox{AV}_{\mbox{\footnotesize box}}\\
\mbox{SE}_{\mbox{\footnotesize av}} & = \mbox{SD}_{\mbox{\footnotesize box}}\,\big{/}\sqrt{n}\\[10pt]
%
\mbox{EV}_{\mbox{\footnotesize \%}} & = \mbox{AV}_{\mbox{\footnotesize box}}\\
\mbox{SE}_{\mbox{\footnotesize \%}} & = \mbox{SD}_{\mbox{\footnotesize box}}\,\big{/}\sqrt{n}\\[10pt]
%
\noalign{Shortcut formula (if there are only two different kinds of numbers in the box):}
\mbox{SD}_{\mbox{\footnotesize box}} &=
  \left(\mbox{big \#} - \mbox{small \#}\right)\,\sqrt{
  \left(\begin{array}{c}\mbox{fraction of}\\ \mbox{tickets with the}\\ \mbox{big number}\end{array}\right)
  \left(\begin{array}{c}\mbox{fraction of}\\\mbox{tickets with the}\\\mbox{small number}\end{array}\right)}\\[10pt]
%
\mbox{95\% confidence interval} &= \mbox{observed} \;\pm\; 2\,\mbox{SE}\\[10pt]
z &= \frac{\mbox{observed} - \mbox{expected}}{\mbox{standard error}}
\end{align*}
\vfill
\eject

%%%%%%%%%%%%%%%%%%%%%%%%%%%%%%%%%%%%%%%%%%%%%%%%%%%%%%%%%%%%%%%%%%%%%%%%%%%

\href{http://www.su.edu}{\includegraphics[height=1.75cm]{sulogo.eps}}
\vspace{-1.69cm}

{{\ }\hfill\small
\begin{tabular}{cl}
& Math 207\\
& Introduction to Statistics\\
\end{tabular}
}
\begin{center}
\textbf{\large  Exam IV}
\end{center}
%No books, notes or electronic devices may be used on this exam.
\medskip

1. (25 points) A survey is conducted of 100 freshmen at public universities.
Among the surveyed  students, 80\% (or 4/5) work at least 5 hours per week for pay.

\hspace{20pt} a) Find a 95\% confidence interval
on the  percent of freshmen at public universities who work at least 5 hours per week for pay.
\vspace{2.3in}

\hspace{20pt} b) True or false and explain:  There is a 95\% chance that the
percent among \textit{all} freshmen at public universities is in the range you found in a).
\vspace{1.5in}

\hspace{20pt} c) True or false and explain:  There is a 95\% chance that
if a second sample of 100 freshmen were taken,
the sample percent would be in the range you found in a).
\vspace{1.5in}

\hspace{20pt} d) True or false and explain:  There is a 95\% chance that 
 among all 18--20 year olds, the percent who work at least 5 hours per week for pay 
is in the range you found in a).
\vspace{1in}
\vfill
\eject
{\ }



2. (10 points)
Would taking the average of 25 measurements decrease
the likely size of the chance error by a factor of 5, 10 or 25? Justify your answer.
\vspace{1.5in}

3. (15 points) A surveyor is measuring the distances between five points $A$, $B$, $C$,
$D$, and $E$ along a straight line.  She finds that each of the four distances measures one 
mile, give or take an inch or so.  The distance from $A$ to $E$ is about four miles; 
but the estimate is likely to be off by around:  
4~inches, 2~inches, 1~inch, 1/2~inch or 1/4~inch. Justify your answer.\vspace{-5pt}

\[\xymatrix{
A\ar@{-}[r] & B\ar@{-}[r] & C\ar@{-}[r] & D\ar@{-}[r] & E}\]
\vspace{2in}


4.  (20 points) The Zorro News Organization reports that  85\%
of freshmen at public universities  work 5 hours a week for pay.
If that is report is accurate, approximately what is the chance of getting a sample
percent as low as or lower than 80\% (as we did in \#1)?  Hint:  calculate $z$ then a $p$-value (area under the normal curve).
\vfill
\eject


5. (20 points) Five hundred draws are made at random from the box
\begin{center}
\begin{pspicture}(0,0.1)(4,0.3)
\psline(0,0.5)(0,0)(4,0)(4,0.5)
\rput[l](0.15,0.3){20,000\ \framebox{0}\,s}
\rput[l](2.15,0.3){60,000\ \framebox{1}\,s}
\end{pspicture}
\end{center}
Note that $\sqrt{0.25\times 0.75}\,/\sqrt{500} \approx 0.02$. True or false?  Explain your conclusions.
\smallskip

\begin{itemize}
\item[i.] The expected value for the percentage of 1s among the draws is exactly 75\%.\\[1in]
\item[ii.] The expected value for the percentage of 1s among the draws is around 75\%, give
or take 2\% or so.\\[1in]
\item[iii.] The percentage of 1s among the draws will be around 75\%, give or take 2\% or so.\\[1in]
\item[iv.] The percentage of 1s among the draws will be exactly 75\%.\\[1in]
\item[v.] The percentage of 1s in the box is exactly 75\%.\\[1in]
\item[vi.] The percentage of 1s in the box is around 75\%, give or take 2\% or so.
\end{itemize}

\vfill
\eject

6. (10 points) Find the specified area under the normal curve.  Write down the \texttt{R} command that you use.

\hspace{10pt} a) $z>1.35$ 
\bigskip
\bigskip\bigskip\bigskip

\hspace{10pt} b) $z<-2$
\bigskip
\bigskip\bigskip\bigskip

\hspace{10pt} c) $-1.5 < z < 3.5$
\vfill



\begin{center}
\textbf{\large Formula Sheet}
\end{center}\vspace{-15pt}
\begin{align*}
\mbox{mean} &=\frac{1}{n}\left(x_1 + \cdots + x_n\right)\\[10pt]
\mbox{SD}   &=\sqrt{\frac{1}{n}\left( (x_1 - \mbox{mean})^2 + \cdots + (x_n - \mbox{mean})^2\right)}\\[10pt]
\mbox{sd}   &=\sqrt{\frac{1}{n-1}\left( (x_1 - \mbox{mean})^2 + \cdots + (x_n - \mbox{mean})^2\right)}\\[10pt]
%
\mbox{EV}_{\mbox{\footnotesize sum}} & = n\,\mbox{AV}_{\mbox{\footnotesize box}}\\
\mbox{SE}_{\mbox{\footnotesize sum}} & = \sqrt{n}\;\;\mbox{SD}_{\mbox{\footnotesize box}}\\[10pt]
%
\mbox{EV}_{\mbox{\footnotesize av}} & = \mbox{AV}_{\mbox{\footnotesize box}}\\
\mbox{SE}_{\mbox{\footnotesize av}} & = \mbox{SD}_{\mbox{\footnotesize box}}\,\big{/}\sqrt{n}\\[10pt]
%
%\mbox{EV}_{\mbox{\footnotesize \%}} & = \mbox{AV}_{\mbox{\footnotesize box}}\\
%\mbox{SE}_{\mbox{\footnotesize \%}} & = \mbox{SD}_{\mbox{\footnotesize box}}\,\big{/}\sqrt{n}\\[10pt]
%
\noalign{Shortcut formula (if there are only two different kinds of numbers in the box):}
\mbox{SD}_{\mbox{\footnotesize box}} &=
  \left(\mbox{big \#} - \mbox{small \#}\right)\,\sqrt{
  \left(\begin{array}{c}\mbox{fraction of}\\ \mbox{tickets with the}\\ \mbox{big number}\end{array}\right)
  \left(\begin{array}{c}\mbox{fraction of}\\\mbox{tickets with the}\\\mbox{small number}\end{array}\right)}\\[10pt]
%
\mbox{95\% confidence interval} &= \mbox{observed} \;\pm\; 2\,\mbox{SE}\\[10pt]
z &= \frac{\mbox{observed} - \mbox{expected}}{\mbox{standard error}}
\end{align*}

\eject


6. (10 points) Find the specified area under the normal curve.  Write down the \texttt{R} command that you use.

\hspace{10pt} a) $z<2.15$ 
\bigskip
\bigskip\bigskip\bigskip

\hspace{10pt} b) $z>-1.75$
\bigskip
\bigskip\bigskip\bigskip

\hspace{10pt} c) $-1.25 < z < 2.5$
\vfill



\begin{center}
\textbf{\large Formula Sheet}
\end{center}


\begin{align*}
\mbox{mean} &=\frac{1}{n}\left(x_1 + \cdots + x_n\right)\\[10pt]
\mbox{SD}   &=\sqrt{\frac{1}{n}\left( (x_1 - \mbox{mean})^2 + \cdots + (x_n - \mbox{mean})^2\right)}\\[10pt]
\mbox{sd}   &=\sqrt{\frac{1}{n-1}\left( (x_1 - \mbox{mean})^2 + \cdots + (x_n - \mbox{mean})^2\right)}\\[10pt]
%
\mbox{EV}_{\mbox{\footnotesize sum}} & = n\,\mbox{AV}_{\mbox{\footnotesize box}}\\
\mbox{SE}_{\mbox{\footnotesize sum}} & = \sqrt{n}\;\;\mbox{SD}_{\mbox{\footnotesize box}}\\[10pt]
%
\mbox{EV}_{\mbox{\footnotesize av}} & = \mbox{AV}_{\mbox{\footnotesize box}}\\
\mbox{SE}_{\mbox{\footnotesize av}} & = \mbox{SD}_{\mbox{\footnotesize box}}\,\big{/}\sqrt{n}\\[10pt]
%
\mbox{EV}_{\mbox{\footnotesize \%}} & = \mbox{AV}_{\mbox{\footnotesize box}}\\
\mbox{SE}_{\mbox{\footnotesize \%}} & = \mbox{SD}_{\mbox{\footnotesize box}}\,\big{/}\sqrt{n}\\[10pt]
%
\noalign{Shortcut formula (if there are only two different kinds of numbers in the box):}
\mbox{SD}_{\mbox{\footnotesize box}} &=
  \left(\mbox{big \#} - \mbox{small \#}\right)\,\sqrt{
  \left(\begin{array}{c}\mbox{fraction of}\\ \mbox{tickets with the}\\ \mbox{big number}\end{array}\right)
  \left(\begin{array}{c}\mbox{fraction of}\\\mbox{tickets with the}\\\mbox{small number}\end{array}\right)}\\[10pt]
%
\mbox{95\% confidence interval} &= \mbox{observed} \;\pm\; 2\,\mbox{SE}\\[10pt]
z &= \frac{\mbox{observed} - \mbox{expected}}{\mbox{standard error}}
\end{align*}


\vfill
\eject

\newcommand{\threebox}[3]{\begin{pspicture}(0.2,0.2)(2.5,1.1)
\psframe(0,0)(0.7,0.7)\rput(0.35,0.35){#1}
\psframe(1,0)(1.7,0.7)\rput(1.35,0.35){#2}
\psframe(2,0)(2.7,0.7)\rput(2.35,0.35){#3}
\psline(-0.3,0.9)(-0.3,-0.2)(3,-0.2)(3,0.9)(-0.3,0.9)
\end{pspicture}}


5. (25 points) One hundred draws will be made at random with replacement
from the box shown below.
\begin{center}
\threebox{0}{0}{1}
\end{center}
\bigskip
Note that $\mbox{SD}_{\scriptsize \mbox{box}}=\sqrt{2}\,/\,3 = 0.47$.
\bigskip


\hspace{20pt} a) 
The number of 1s among the draws will be around \underline{\hspace{45pt}} 
give or take \underline{\hspace{45pt}}  or so.
\vspace{2.5in}

\hspace{20pt} b) 
The percentage of 1s among the draws will be around \underline{\hspace{45pt}} 
give or take \underline{\hspace{45pt}}  or so.
\vspace{2.5in}

\hspace{20pt} c) Roughly what is the probability that the percentage of 1s will 
be over 38\%.




\end{document}

